\section{Background, Related Work, and the Addressed Gap}
\label{sec:background}

\subsection{Modular self-reconfigurable robots}

Modular self-reconfigurable robot systems combine repeated hardware units with
mechanisms for changing connectivity. Their potential benefits include functional
reuse, task-dependent morphology, redundancy, repair, and scalable manufacturing;
their enduring challenges include mechanical connection, configuration planning,
control, communication, and the combinatorial growth of possible morphologies
\citep{yim2007modular}. A configuration is naturally represented by a graph whose
vertices denote modules and whose edges denote connections, but a useful physical
description must also preserve connection faces, orientation, joint state, spatial
frames, and actuation.

SMORES and \smoresep{} are mobile-style modular robots whose left and right faces also
serve as wheels. The original SMORES design demonstrated locomotion and
reconfiguration \citep{davey2012smores}. \smoresep{} adds electro-permanent-magnet
connectors. Each module has four active rotational degrees of freedom---LEFT, RIGHT,
PAN, and TILT---and four mating faces---LEFT, RIGHT, TOP, and BOTTOM
\citep{liu2020parallel,liu2023smoresep,modlabSmoresEp}. LEFT and RIGHT provide
differential-drive locomotion; LEFT, RIGHT, and PAN rotate continuously, while TILT is
bounded.

The EP-Face was designed as a planar, genderless, switchable connector. Published
characterization reports high holding force, rapid state switching, and tolerance to
position error \citep{tosun2016epface}. These properties motivate a simulation model
with explicit connector frames, capture conditions, release timing, and post-capture
constraints. They also illustrate why a bare URDF joint tree is insufficient: a
connector's compatibility, allowable roll orientations, capture region, and release
policy are not standard URDF concepts.

\subsection{Planning topology versus executing motion}

Liu, Whitzer, and Yim formulate reconfiguration as a sequence of connector changes and
demonstrate a distributed planning approach \citep{liu2019distributed}. Their
Driver-to-Snake example gives an initial and target configuration and four connector
replacement actions. This is valuable topology-level ground truth, but it does not
fully specify motor torques, contact coefficients, obstacle-clearance routes, or a
time-parameterized trajectory for the entire Driver-to-Snake sequence. Later SMORES
self-assembly work develops navigation, pose adjustment, and straight-approach control
and demonstrates hardware behaviors \citep{liu2020parallel,liu2023smoresep}.

The distinction between \emph{what connections should change} and \emph{how bodies
should move} is central to \modsim{}. A \code{ReconfigurationPlan} describes connector
replacements. A physical scenario supplies platform mechanics, control parameters, and
routes. The semantic runtime remains responsible for validating and committing the
connection. This separation allows future planners to replace authored routes without
changing the docking or event model.

\subsection{Physics engines and robot-description formats}

Modern simulators already solve difficult numerical problems: articulated dynamics,
contact, equality constraints, actuators, sensors, and visualization. \mujoco{} is
designed for efficient model-based control and represents multijoint dynamics in
generalized coordinates \citep{todorov2012mujoco}. Its native MJCF language can express
rich dynamic models, and its compiler can also import URDF
\citep{mujocoModelingDocs,urdfFormat}. Isaac Sim and related USD-based workflows offer
another path to high-fidelity scene composition and rendering. Gazebo/SDF, Webots, and
CoppeliaSim provide still other combinations of physics, scene, and control tooling.

These engines intentionally operate below the domain model proposed here. A weld
constraint can hold two bodies together, but it does not by itself establish:

\begin{itemize}
  \item that two surfaces are semantically compatible;
  \item that their approach axes and roll orientations satisfy a connector contract;
  \item that a runtime command or auto-latch policy allows connection;
  \item that the logical module graph should gain a particular stable edge; or
  \item that downstream algorithms should receive a canonical dock event.
\end{itemize}

URDF has a complementary limitation. It is widely supported as a mechanical
description of links, joints, inertial data, and geometry, but it does not standardize
modular-robot docking semantics. MJCF can encode engine-specific contact and actuator
details, but using MJCF alone would make the portable semantic package engine-specific.
\modsim{} therefore treats URDF/MJCF as mechanical assets referenced by a separate
semantic package.

\subsection{Why an engine-independent semantic layer is needed}

The addressed gap lies between a robot asset and a modular-robot application. The asset
describes what can move and collide. The application needs stable identities,
compatibility rules, topology, lifecycle events, validation, and generated views.
Embedding those concepts directly in a backend creates several failure modes:

\begin{enumerate}
  \item \textbf{Duplicated truth.} A graph, a scene graph, and a controller each carry a
        different list of active docks.
  \item \textbf{Optimistic commit.} The application records a dock even though the
        engine refuses or cannot allocate a physical constraint.
  \item \textbf{Frame drift.} A connector pose authored in one tool is recomposed
        differently by another tool, so visual alignment and physical acceptance
        disagree.
  \item \textbf{Non-portable scenes.} Connector semantics are hidden in simulator
        scripts and cannot be validated before that simulator is installed.
  \item \textbf{Uninspectable failures.} A boolean docking result omits which criterion
        failed and by how much.
  \item \textbf{Unsafe extensibility.} A data package contains arbitrary executable
        plugin paths rather than stable builder or backend identifiers.
\end{enumerate}

\subsection{Positioning relative to existing simulators}

\Cref{tab:positioning} summarizes the intended division of responsibility. The table
does not claim that general simulators cannot be extended to provide these features;
rather, it identifies the contract \modsim{} adds above their native physical model.

\begin{table}[htbp]
\centering
\caption{Conceptual positioning of \modsim{} relative to a physics engine.}
\label{tab:positioning}
\begin{tabularx}{\textwidth}{L{0.25\textwidth}YY}
\toprule
Concern & Physics engine or robot asset & \modsim{} semantic layer \\
\midrule
Rigid-body state & Integrates generalized coordinates, contact, and constraints &
Ingests observations into stable module/link/connector identities \\
Mechanical description & URDF, MJCF, or a future USD asset &
References assets and adds portable module semantics \\
Dock validity & Can test distance or create a raw constraint &
Defines compatibility, acceptance, policy, lifecycle, and failure detail \\
Dynamic connection & Activates/deactivates an equality or joint &
Uses two-phase commit before mutating logical topology \\
Configuration & Scene/body graph & Canonical module multigraph and assemblies \\
History & Engine-dependent logging & Ordered semantic event log and metrics \\
Visualization & Native 3-D scene & Derived, immutable topology and event views \\
Authoring & Engine-specific editor or XML & Schema-aware Robot Pack Studio \\
\bottomrule
\end{tabularx}
\end{table}

The resulting architecture is intentionally compositional. \modsim{} benefits from a
mature solver instead of reimplementing one; the solver benefits from a structured
domain layer that can be tested independently. A future C++ execution core, browser
client, or second simulator can reuse the same semantic contracts.

